\documentclass{article}
\usepackage[margin=1.15in]{geometry}
\usepackage{xcolor}
\usepackage{parskip}
\usepackage{fancyhdr}
\usepackage{array}
\usepackage{booktabs}
\usepackage{enumitem}
\usepackage{amsmath}
\usepackage{tabularx}
\usepackage{listings}

\pagestyle{fancy}
\fancyhf{}
\cfoot{\thepage}
\rhead{Lecture 10}
\lhead{MA 214: Applied Statistics (Spring 2026)}
\renewcommand{\headrulewidth}{.4mm}

\newcommand{\blankline}[1]{\underline{\hspace{#1}}}

\lstset{
  basicstyle=\ttfamily\small,
  columns=fullflexible,
  frame=single,
  framerule=0.4pt,
  breaklines=true,
  showstringspaces=false
}

\begin{document}

\noindent\textbf{Name:}\blankline{2.25in}\hfill\textbf{BUID:}\blankline{1.55in}

\medskip

\noindent\textbf{Name:}\blankline{2.25in}\hfill\textbf{BUID:}\blankline{1.55in}\newline

\begin{center}
 \Large In-Class Worksheet: Inference with Mathematical Models
\end{center}

\subsection*{Goals}
\noindent\textbf{(1)} compute exact and approximate probabilities, \textbf{(2)} turn a standard error into an interval and a p-value, \textbf{(3)} read and debug the R code that does it, \textbf{(4)} tell paired data from two samples, \textbf{(5)} know when a formula is not allowed.

\subsection*{Part 1: Exact, approximate, simulated}
A biologist traps $n=20$ rodents and finds $14$ males. She wants to test $H_0: p = 0.5$.

\begin{enumerate}[label={\bfseries (\Alph*)},itemsep=.8em]
\item Which model gives the \emph{exact} null distribution of the number of males? Write it down, with its parameters. \\[2.5em]
\item She instead simulates $1{,}000$ datasets under $H_0$ and counts. Would her p-value equal the exact one? Explain in one sentence what would change if she re-ran the simulation. \\[3.5em]
\item Check the condition for using a \emph{normal approximation} here ($np \ge 10$ and $n(1-p) \ge 10$, computed under $H_0$). Is it satisfied? \\[2.5em]
\end{enumerate}

\subsection*{Part 2: From a standard error to an interval}
Last week we bootstrapped the mean weight of $n = 348$ kangaroo rats. This week we have the summary statistics instead: $\bar{x} = 43.95$ g and $s = 6.56$ g.

\begin{enumerate}[label={\bfseries (\Alph*)},itemsep=.8em]
\item Compute the standard error of the mean. \hfill $\mathrm{SE} = \blankline{1.2in}$ \\[1em]
\item Build the 95\% confidence interval $\bar{x} \pm 1.96\,\mathrm{SE}$: \quad $\big[\;\blankline{.8in}\;,\;\blankline{.8in}\;\big]$ grams. \\[1em]
\item Last week's bootstrap percentile interval was $[43.27,\ 44.63]$. Compare. Why should we \emph{expect} these to agree --- and what would it have meant if they had not? \\[4em]
\item You want an interval half as wide. By what factor must $n$ grow? (Look at where $n$ sits in the formula.) \hfill \blankline{1.2in} \\[1.5em]
\end{enumerate}

\newpage

\subsection*{Part 3: Find the bugs}
A classmate wants a 95\% confidence interval for the mean weight. The vector \texttt{w} holds the $348$ weights. Their code runs without an error message --- and is wrong in \textbf{three} separate ways.

\begin{lstlisting}[language=R]
xbar <- mean(w)
se   <- sd(w)                      # standard error

xbar + c(-1, 1) * 1.96 * se        # 95% CI

# and to test H0: mu = 40
z <- (xbar - 40) / se
pnorm(-abs(z))                     # p-value
\end{lstlisting}

\begin{enumerate}[label={\bfseries (\Alph*)},itemsep=.8em]
\item Bug 1 is on the \texttt{se} line. What is missing, and is their interval too wide or too narrow? \\[3.5em]
\item Bug 2: $s$ was estimated from the data. Which critical value should they have used instead of $1.96$, and would it make the interval wider or narrower? \\[3.5em]
\item Bug 3 is on the last line. What kind of p-value does \texttt{pnorm(-abs(z))} give, and what should it be for a two-sided test? \\[3em]
\item Rewrite the corrected code: \\[5em]
\end{enumerate}

\subsection*{Part 4: Paired, or two samples?}
For each study, circle one and name the statistic you would compute.

\begin{enumerate}[label={\bfseries (\Alph*)},itemsep=1em]
\item Cholesterol of 40 patients, measured before and after a diet. \hfill \textbf{paired / two-sample} \\
Statistic: \blankline{3in}
\item Weights of the 204 male and 144 female rats trapped in 1999. \hfill \textbf{paired / two-sample} \\
Statistic: \blankline{3in}
\item Prices of the same 73 textbooks at two different shops. \hfill \textbf{paired / two-sample} \\
Statistic: \blankline{3in}
\item In one sentence: what is the \emph{test} for whether data can be paired? \\[3em]
\end{enumerate}

\subsection*{Part 5: When is a formula not allowed?}
For each goal, say whether today's formulas apply. If they do not, say what you would do instead --- and be honest about what that costs.

\begin{enumerate}[label={\bfseries (\Alph*)},itemsep=.8em]
\item A confidence interval for the \emph{90th percentile} of rat weight, from all 348 rats. \\[3em]
\item A p-value for whether a coin is fair, from 200 flips. \\[3em]
\item A confidence interval for the mean weight of $n = 6$ rats trapped at a new site. \\
\emph{Careful.} You have three tools now. Say what each one would require of you here --- and whether bootstrapping is really a way out. (Look back at the $n=5$ picture from Lecture 9.) \\[5em]
\end{enumerate}

\end{document}
